\section{结论与讨论}

Depth Anything 3 表明，一个简单的Transformer，在深度与射线目标的教师-学生监督下进行训练，无需复杂的架构即可统一任意视角的几何信息。尺度感知的深度、逐像素射线以及自适应跨视角注意力使模型能够继承强大的预训练特征，同时保持轻量且易于扩展。
在所提出的视觉几何基准上，该方法创下了新的位姿估计与重建记录，其大型和紧凑变体均超越了先前的模型，同时同一骨干网络还支撑了高效的前馈式新视角合成模型。

我们将 Depth Anything 3 视为迈向通用3D基础模型的一步。未来的工作可以将其推理能力扩展到动态场景，整合语言与交互线索，并探索更大规模的预训练，以打通几何理解与可操作世界模型之间的闭环。我们希望本文提供的模型与数据集发布、基准测试以及简洁的建模原则能够推动通用3D感知的更广泛研究。

\section*{致谢}
我们感谢周晓伟、彭思达和郭恒凯在本项目开发过程中提供的宝贵讨论。
我们也感谢赵杨在工程上的支持。
预告演示中的输入图像提取自一个公开的YouTube视频\citep{mtsdrones2024}，鸣谢原作者。
